\documentclass[11pt]{article}
\usepackage[margin=1in]{geometry}
\usepackage{amsmath,amssymb,amsthm,mathtools}
\usepackage[colorlinks=true,urlcolor=blue,linkcolor=blue]{hyperref}
\newtheorem{lemma}{Lemma}
\newtheorem{proposition}[lemma]{Proposition}
\newtheorem{corollary}[lemma]{Corollary}
\newcommand{\lgx}{\log_2}
\newcommand{\MS}{\operatorname{MS}}
\newcommand{\OPT}{\operatorname{OPT}}
\newcommand{\phiu}{\Phi_{\mathrm{unif}}}
\title{Multisplay versus Tango squared:\\emission order and potential obstructions}
\author{Research notes, 2 October 2026}
\date{}
\begin{document}
\maketitle

\noindent\textbf{Status.} These are proved elementary partial statements and
explicit obstructions to particular amortization schemes. They neither prove
nor disprove a third-iterated-logarithm competitive ratio for standard
multisplay. The bit-count statements refer to precisely specified words;
they do not identify those words with the searches performed by multisplay.

\section{Changing the order of a fixed set of emissions}

Fix a full binary hierarchy $B$ on the ordered depth labels
$0,\ldots,H$. An internal interval $I=[a,b]$ has a left and right child.
For each original request $x_t$, let
\[
 E_t=\{d_{t,1}<\cdots<d_{t,k_t}=H\}
\]
be a nonempty set of distinct emitted depths. A write of $d\in E_t\cap I$
goes to the fixed indexed bit $(I,\operatorname{anc}_a(x_t))$ and writes
the side of $I$ containing $d$. All indexed bits start at R.
All three schedules below use the same $E_t$, owners, and original-request
order. Within a request, the schedules are
\[
 \uparrow=(d_1,\ldots,d_k),\qquad
 \downarrow=(d_k,\ldots,d_1),\qquad
 \operatorname{rev}H=(d_{k-1},\ldots,d_1,H).
\]
Their total initialized bit-change counts are $F^\uparrow$, $F^\downarrow$,
and $F^{\operatorname{rev}H}$. Call an interval mixed for request $t$
if $E_t$ meets both its children. Let $M_t$ count all mixed intervals and
$M_t^{\mathrm{nt}}$ count the mixed intervals not containing $H$.

\begin{lemma}[Order perturbation]
For the fixed owner-indexed words above,
\[
 |F^\downarrow-F^\uparrow|\le 2\sum_t M_t
       =2\sum_t(k_t-1),
\]
and
\[
 |F^{\operatorname{rev}H}-F^\uparrow|
   \le 2\sum_t M_t^{\mathrm{nt}}
   \le 2\sum_t\max\{k_t-2,0\}.
\]
Moreover, every terminal indexed bit, meaning $H\in I$, has exactly the
same change count under $\uparrow$ and $\operatorname{rev}H$.
\end{lemma}
\begin{proof}
Consider one fixed indexed bit and divide its word into nonempty request
blocks. A pure block is unchanged. A mixed increasing block is
$L^pR^q$ with $p,q>0$. Full reversal replaces it by $R^qL^p$.
Both have one internal change. Only the incoming and outgoing block
boundaries can change their contribution, by at most one each. Replacing
the mixed blocks one at a time proves the first inequality. This argument
includes the initial R as the incoming boundary of the first block;
there is no outgoing boundary after the final block.

For a nonterminal interval, keeping $H$ last makes no difference from
full reversal on its filtered word, so the same argument applies.
For a terminal interval, every nonempty block in this indexed row contains
the completion $H$. Indeed, any request to this owner writes its completion
to this interval. Thus both schedules end every block in R and the row
begins in R. An all-R block contributes zero. A block with any L contributes
exactly two changes: under $\uparrow$ its form is $L^pR^q$; under
$\operatorname{rev}H$ it is $R^{q-1}L^pR$, with a possibly empty initial
R run. The two contributions are identical, block by block.

The mixed intervals are exactly the branching nodes in the minimal
subtree of $B$ spanning $E_t$. Suppressing unary vertices yields a full
binary tree on $k_t$ leaves, so $M_t=k_t-1$. If $k_t\ge2$, at least one
branching node is an ancestor of $H$. Therefore
$M_t^{\mathrm{nt}}\le k_t-2$. If $k_t=1$, both quantities vanish.
\end{proof}

For the canonical Tango-squared emissions,
$\sum_t(k_t-1)=P+I_0$ with $I_0\le n-1$ in the complete-universe
convention. The lemma therefore preserves an $O(m+P+D+n)$ initialized
secondary budget. In particular, on a normalized request word with at
most one noncompletion emission per request,
$F^{\operatorname{rev}H}=F^\uparrow$ exactly.

\paragraph{Scope.}
Standard multisplay also splays reference-ancestor boundary keys, and its
internal-node requests have a final switch convention. Neither operation
is identified with a completion at depth $H$ by this lemma. A physical
implementation and a word-identification theorem remain necessary.

\section{The global range of the uniform multisplay potential}

Here use the standard node-storing reference: a perfect binary tree of
height $h$, with $n=2^{h+1}-1$ nodes. Every internal vertex selects one
preferred child. Its preferred paths are represented by auxiliary splay
trees. For $v$ in auxiliary $A$, let $s_A(v)$ count the nodes in its splay
subtree, excluding marked child components. Set
\[
 \phiu=\sum_A\sum_{v\in A}\lgx s_A(v).
\]

\begin{lemma}[Uniform potential has linear global range]
For every valid preferred-path decomposition and every choice of auxiliary
BST shapes, $0\le\phiu=O(n)$. The multiset of auxiliary sizes is independent
of the preferred-child choices.
\end{lemma}
\begin{proof}
The root path has $h+1$ nodes. At each depth $a\ge1$, precisely
$2^{a-1}$ vertices are nonpreferred children, hence are path heads.
Every such path has $h-a+1$ nodes. Thus the size multiset is fixed.
An auxiliary of size $\ell$ contributes at most $\ell\lgx\ell$.
Consequently
\[
 \phiu\le(h+1)\lgx(h+1)+
       \sum_{\ell=1}^{h}2^{h-\ell}\ell\lgx\ell=O(2^h)=O(n),
\]
because $\sum_{\ell\ge1}2^{-\ell}\ell\lgx\ell$ converges.
\end{proof}

Thus this potential needs only an $O(n)$ initialization allowance.
This observation does not improve its per-request amortized estimate.
It also means that adding any term $\sum_A g(|A|)$ is just a constant
normalization in this perfect-reference setting.

\section{The exact marking increment of a switch}

Consider a left-to-right switch at reference vertex $y$, after its at most
three splays have exposed $y$ as auxiliary root and, when present,
$x=\operatorname{refLeftParent}(y)$ and
$z=\operatorname{refRightParent}(y)$ as its left and right children.
Let $a$ and $b$ be the numbers of strict ancestors of $y$ in its current
preferred path lying respectively to its left and right in key order.
Both old and incoming suffixes have the same size $t$ in a perfect
reference tree. For $s>0$ define $f_t(s)=\lgx(1+t/s)$, and define
$f_t(0)=0$ to express the absence of that boundary node.

\begin{lemma}[Marking cross-ratio]
The change in $\phiu$ due solely to marking the old suffix and unmarking
the incoming suffix, after the exposing splays, is
\[
 J_{L\to R}=f_t(b)-f_t(a).
\]
For a right-to-left switch it is $J_{R\to L}=f_t(a)-f_t(b)$.
\end{lemma}
\begin{proof}
If $x$ exists, its size decreases from $a+t$ to $a$. If $z$ exists,
its size increases from $b$ to $b+t$. The size of $y$ returns to its
original value because the exchanged suffixes have equal sizes.
No other node's auxiliary-subtree size changes. Summing these logarithmic
changes proves the formula. Reflecting key order proves the other case.
\end{proof}

For example $a=t,b=1$ gives
$J_{L\to R}=\lgx(t+1)-1$. Once the frame is exposed, the physical
marking work is constant even though its potential increase is
$\Theta(\lgx t)$. A normalizer depending only on auxiliary sizes cannot
alter this increment.

\section{A full-request oscillation on two leaves}

The same issue occurs for whole requests, not just isolated mark changes.
Take $h\ge2$. Let $r$ be the reference root, $x$ its globally leftmost
leaf, and $y$ the leftmost leaf of its right reference subtree. The two
root-child suffix paths each contain $h$ keys. Let $p_x,p_y$ be the
reference parents of the leaves; let $A_L,A_R$ hold the other $h-2$ keys
on the respective suffix paths. Keep all hanging off-path auxiliaries
opaque. Write $v(L,R)$ for an auxiliary node, omitting marked children.
The following two states form a period for requests $x,y$:
\[
\begin{array}{c|c|c}
 &\text{root auxiliary}&\text{other displayed auxiliary}\\ \hline
 X\ (\text{after }x)&
 x(\varnothing,p_x(\varnothing,r(A_L,\varnothing)))&
 y(\varnothing,p_y(\varnothing,A_R))\\[2pt]
 Y\ (\text{after }y)&
 y(r,p_y(\varnothing,A_R))&p_x(x,A_L).
\end{array}
\]
All omitted attachments stay in their actual ordered gaps.

\begin{proposition}
Standard bottom-up multisplay changes $X$ to $Y$ on request $y$, and
$Y$ to $X$ on request $x$, using three rotations per request and constant
navigation and marking work. Yet
\[
 \phiu(X)-\phiu(Y)=\lgx h+\lgx(h-1).
\]
\end{proposition}
\begin{proof}
In $X$, the route to $y$ goes through $x,p_x,r$ and the marked right
suffix. Splaying $r$ uses one zig-zig, leaving the old left suffix
$p_x(x,A_L)$. Mark it, unmark the right suffix, and splay $y$ by one
zig. Its two reference-ancestor boundaries are already $r$ and $p_y$,
so its final leaf switch needs no further rotations.

In $Y$, the route to $x$ goes through $y,r,p_x$. Splay $r$ by one zig,
mark the right suffix, and unmark the left one. Splaying $x$ then uses
one zig-zig; its sole reference-ancestor boundary $p_x$ is already its
right child. This gives precisely $X$.

The two active-root contributions are $\lgx(h+1)$ and cancel. The
contributions of $p_x$ and $p_y$ are respectively $\lgx h$ and
$\lgx(h-1)$ in both states. The opaque subtrees also cancel. The extra
contributions in $X$ are $\lgx(h-1)$ at $r$ and $\lgx h$ at the dormant
root $y$; their counterparts $r,x$ in $Y$ have size one.
\end{proof}

After a finite warmup, the canonical leaf-reference histories of these
two keys have one primary change per request, at their fixed LCA.
Their depth emissions are its fixed depth and $H$. Thus their initialized
secondary count is $O(1)$ per request after warmup. In particular, placing
the two keys in different halves of the complete leaf reference gives
the emissions $(0,H)$ and exactly two secondary changes per request.
The node-storing reference used for multisplay and the leaf reference
used for this count are explicitly different descriptions; only this
two-key history claim is used here.

It follows that the proposed \emph{pointwise} estimate
\[
 \operatorname{cost}_{\MS}(t)+\Delta\phiu(t)
 \le C(1+\lgx\lgx(H+2))(1+p_t+f_t)
\]
cannot hold with a universal $C$. The actual algorithm is cheap on this
example: the negative potential jump on the opposite request cancels the
positive one. This is an obstruction to this potential inequality, not
an algorithmic lower bound.

\section{A preference-only correction cannot absorb marking separately}

\begin{proposition}[Nonzero circulation of marking increments]
There is no state function $G$ of the preferred-edge assignment alone
whose change cancels $J$ on every switch. More strongly, raw marking
increments plus a telescoping such $G$ cannot be charged
$O((m+P+F)(1+\lgx\lgx(H+2)))$ on every repeated request period.
\end{proposition}
\begin{proof}
In a perfect node reference of height $h\ge2$, let $u$ be the root and
$v$ its left child. Start with both preferring left. Perform the cycle
\[
 (v:L\to R),\quad(u:L\to R),\quad
 (v:R\to L),\quad(u:R\to L).
\]
For the first switch $a=0,b=1,t=h-1$, so $J=\lgx h$.
Both root switches have no strict reference ancestors in their auxiliary,
so their marking increments are zero. During the third switch $v$ is
the head of a detached path, so its marking increment is also zero.
The preference assignment returns to its initial state, but
\[
 \sum_{\mathrm{cycle}}J=\lgx h.
\]
Therefore no such exact state correction exists.

This is a realizable query schedule, not merely a formal switch cycle.
Choose leaves $a,b,c$ respectively in $v$'s left subtree, $v$'s right
subtree, and $u$'s right subtree. Fix all lower preferences toward these
leaves and repeat the requests $(b,c,a)$ starting after $a$. Their wrong
switches are $(v)$, $(u)$, and $(v,u)$ in bottom-up order. Final switches
at leaves change no preferred edges or root markings.

For example choose the leftmost leaf in each of the three regions.
The node reference has keys $1,\ldots,2^{h+1}-1$; these three queried
keys have zero-based ranks $0,2^{h-1},2^h$. Embed those ranks in the
complete leaf reference with $2^{h+1}$ leaves, height $H=h+1$, and
leave its final leaf unused. The steady emitted sets are
$\{1,H\},\{0,H\},\{0,1,H\}$. There are four primary flips
per period. Only the secondary separation of $\{0,1\}$ and the terminal
separation from $H$ can alternate, giving a bounded secondary count
(eight per period for the midpoint hierarchy with $H\ge3$).
Every $G$ telescopes to zero on a period, while the marking charge remains
$\lgx h$. Since $\lgx h/(1+\lgx\lgx(h+2))$ is unbounded, the stronger
claim follows. Repeating sufficiently many periods absorbs any allowed
one-time $O(n)$ or $O(n(1+\lgx\lgx(H+2)))$ initialization term.
Auxiliary shapes can change across the period; this does
not change $J$, which depends only on the reference-path configuration.
\end{proof}

The multisplay rotations must therefore participate in any successful
amortization. Charging all splays by a nonnegative bound and then seeking
a separate preference-only correction for the marking bill loses essential
cancellation. The proposition permits a correction depending also on
physical shapes, or a joint inequality for rotations and markings.

\section{The missing upper statement}

The Tango-squared count theorems give
$m+P+D=O(\OPT)$ and $F=O(D+P+n)$. A sufficient new upper theorem would be
\[
 \MS(X)=O\bigl((m+P(X)+F(X)+n)
                    (1+\lgx\lgx(H+2))\bigr),
 \qquad H=\Theta(\lgx n).
\]
One possible proof would give a single global potential coupling the
actual multisplay forest to the fixed-owner depth histories, of bounded
initial deficit, and an amortized inequality for complete requests.
The companion note \path{endpoint-embedding.tex} gives a closer lower-side
target by representing each standard multisplay request with one of two
adjacent copies in a full leaf reference. Let $S$ be its actual switch count
and let $F_{\MS}$ count the fixed-owner writes in actual bottom-up ancestor
switch order followed by the final switch at the requested node. That final
switch has the largest depth in its block; it is not a global depth-$H$
completion. The endpoint theorem and the analogous order lemma supply
the order lemma above supply
\[
 m+S+F_{\MS}=O(\OPT(X)+n).
\]
The corresponding unproved physical upper target is
\[
 \MS(X)=O\bigl((m+S+F_{\MS}+n)
                    (1+\lgx\lgx(H+2))\bigr).
\]
This endpoint count is an analytical statistic; multisplay does not
literally maintain its secondary bits.
The uniform potential and preference-only marking correction above do not
give this theorem. The order lemma removes one purely chronological
obstacle, while preserving the need to identify actual auxiliary boundary
searches, retained histories, and physical split/join work.

\paragraph{Sources and conventions.}
The standard switch and uniform potential are from Wang's dissertation,
\emph{Multi-Splay Trees}, Sections 3.2--3.4.1:
\url{https://www.csd.cmu.edu/sites/default/files/phd-thesis/CMU-CS-06-140.pdf}.
The fixed-owner count conventions and the established Tango-squared
implementation are in the unified manuscript of the neighboring
\texttt{tango\_squared} repository (1 October 2026), in
\path{research/tango-squared.tex}.
The latter upper construction has its own exposed splice; it does not
state this upper bound for standard multisplay.
\end{document}
